\chapter{绪论}
\label{ch:introduction}

\section{课题背景及研究目的与意义}
\subsection{研究背景}

在灾后应急观测、低空巡检和临时通信覆盖等任务中，无人机能够较快到达地面人员难以进入的区域，采集现场图像与传感数据，并为任务区域提供通信支持\cite{intro3,zeng2016uav}。多架无人机通过无线链路组成飞行自组网后，可以根据任务位置调整覆盖范围，在缺少稳定地面基础设施时维持信息交换。部分节点承担局部观测，部分节点同时提供数据中继，地面或边缘中心根据各节点回传的信息形成判断。这一过程将感知、传输和协同决策联系起来，其中任何环节的能力不足都可能影响任务完成。

与位置相对稳定的无线传感网络相比，飞行节点之间的相对速度较高，邻居关系与链路质量变化较快。当前距离目的地较近的中继，可能在后续时隙离开通信范围；当前速率较高的链路，也可能因运动、干扰或带宽变化而失去优势。维护全局拓扑需要持续交换状态，而收集完成的拓扑可能已经过期。因而，无人机路由既需要及时利用新的连接机会，也需要控制信息更新和决策本身的开销。

机载资源进一步限制了网络运行。无人机的飞行、计算和通信均依赖有限电量，任务规划需要为通信保留相应预算。在给定通信预算下，靠近汇聚中心或处于关键位置的节点通常承担更多转发，其电量下降速度可能明显高于外围节点。关键中继退出后，即使其他节点仍有能量，也可能因回传路径断开而无法继续服务。因此，通信性能不仅由当前可达性决定，还与能量消耗在节点之间及任务时间内的分布有关。

传统数据回传通常以图像帧、传感器读数和状态分组为对象，发送内容在路由决策之前已经确定。路由主要回答数据经过哪些节点、使用多大功率等问题。随着机载语义处理能力提高，节点还可以先提取局部观测中的任务信息，将图像转换为类别判断、置信度、简洁陈述或摘要。接收端不必恢复所有原始内容，而可利用这些线索逐步完成判断。任务语义通信由此把优化对象从既定比特流扩展到具有不同价值和表达粒度的消息\cite{semantic-communications-jsac,semantic-communications-survey}。

以灾后道路状况判断为例，一架无人机只能观察入口，另一架能够看到中段积水。单一视角可能形成不完整结论，完整上传各视角又需要较高带宽。如果先提取“入口可见”“存在障碍”等证据，中心便能通过少量信息补充判断。但内容的作用取决于中心已经知道什么：重复描述已确认的区域，通常不如补充盲区或解释冲突有效。短标签成本较低，长摘要保留的依据较多，适宜的粒度会随任务需求、剩余预算和链路条件变化。

\begin{figure}[htbp]
\centering
\begin{tikzpicture}[font=\small,>=Stealth,
 n/.style={draw,rounded corners=2pt,align=center,minimum height=12mm,text width=29mm},
 edge/.style={->,thick},note/.style={align=center,font=\footnotesize}]
\node[n] (a) at (0,2.6) {无人机 A\\局部观测与证据};
\node[n] (b) at (0,0) {无人机 B\\局部观测与证据};
\node[n] (c) at (4.5,1.3) {感知／中继节点\\路由与功率选择};
\node[n] (d) at (9,1.3) {边缘决策中心\\语义融合与判断};
\draw[edge] (a) -- node[note,above,sloped]{时变链路} (c);
\draw[edge] (b) -- node[note,below,sloped]{有限能量} (c);
\draw[edge] (c) -- node[note,above=4pt,text width=14mm,inner sep=0pt]{数据／\\语义消息} (d);
\draw[edge,dashed] (d.south) -- ++(0,-1.2) -| node[note,pos=.25,below]{累计确认与状态反馈} (c.south);
\node[note,draw,dashed,inner sep=5pt,text width=105mm] at (4.5,-2) {网络传输：下一跳与功率\qquad 内容选择：来源与粒度\\交付控制：新发、重传、等待与停止};
\end{tikzpicture}

\bicaption{资源受限无人机网络中的数据回传与语义协同场景}{Data relaying and semantic collaboration in a resource-constrained UAV network}
\label{fig:intro:scenario}
\end{figure}

图~\ref{fig:intro:scenario} 给出了本文关注的协同通信过程。底层网络分配数据的传输路径与功率，语义选择决定哪些内容进入传输，交付控制则根据发送和反馈判断后续动作。三类决策分别作用于通信承载、信息内容和实际接收状态，共同影响有限资源能够支持多少有效任务。

实际链路还会使内容价值与最终收益产生差距。数据发送完成只表示本地尝试结束，消息可能未被完整接收；接收端即使已经获得数据，其确认也可能延迟或丢失。发送端无法立即区分“数据没有到达”和“确认尚未返回”。如果未确认就补发，可能重复传输已有内容；如果持续等待，则可能错过其他证据的发送机会。任务截止时间越短，这种取舍越直接地影响最终结果，也使语义通信需要进一步考虑交付过程。

\subsection{研究问题的提出}

上述场景中的资源限制并非彼此独立。较高功率可能提高当前速率，却加快节点电量下降；较长消息可能补充语义细节，却压缩后续通信窗口；更多重传可能提高单条证据的到达概率，却挤占其他候选及后续任务的资源。本文围绕这些取舍，提出三个相互衔接的研究问题。

（1）动态拓扑下如何协调路由效率与节点能量持续性？地理路由和学习路由分别通过位置规则或交互经验选择路径\cite{gpsr,drlfr,drlrr}，但对连续任务而言，当前链路有利并不意味着可以长期反复使用。按距离选择容易集中使用关键中继，按次数轮询又不能反映链路之间的功率和速率差异。当数据积压增加时，发送过于保守会引起溢出，持续提高功率则可能提前耗尽预算。该问题需要把业务需求、节点历史支出和邻居状态共同纳入局部决策，在未知未来拓扑的条件下联合选择下一跳与发射功率，使当前转发收益与任务所需的服务周期相协调。

（2）在通信预算和链路条件限制下，如何选择具有补充作用的语义证据？局部节点拥有不同视野，同一问题可能产生互补判断，也可能产生共同误判。高置信度并不保证来源可靠，单独作用较强的消息在已有证据之后也可能变得冗余。与此同时，节点和消息粒度的组合形成较多候选，消息长度又通过时延影响送达风险。因而，选择策略不仅需要比较一次发送的收益，还需根据中心信念、历史消息和剩余预算调整整个序列，确定发送来源、表达粒度和停止时机。该问题把第一问中的资源分配进一步落实到通信内容，使有限传输机会优先服务于当前判断。

（3）在确认延迟或丢失时，如何兼顾正确交付与跨任务能耗？消息被选中后，发送端仍无法直接知道接收中心已获得哪些证据。固定补发次数难以适应不同链路和剩余时限，而只使用最新确认又会忽略尚未确认的历史尝试。此外，继续通信应以能否改善任务结果为依据。中心即使收到更多消息，也可能维持原有判断；若仅凭证据数量或置信度上升估计收益，控制器便可能持续发送缺少实际增益的内容。该问题需要区分内容是否新增、是否已经到达以及确认能否及时返回，联合决定新发、重传、等待和停止。能量还应在连续任务间保留状态，避免当前低收益通信占用后续任务的机会。

三问由同一协同过程中的不同矛盾引出。第一问关注传输机会如何持续可用，第二问关注哪些证据值得使用这些机会，第三问关注潜在收益如何在不确定链路上实现。单独提高吞吐量不能说明选中了更有用的内容，选中有用内容也不能保证实际到达。因此，三项研究分别界定动作与约束，再通过共同的任务背景和资源状态建立联系。

\subsection{研究目的与意义}

本文以资源受限无人机网络为对象，研究能够随业务、任务需求和反馈状态调整的通信决策方法。研究目的包括：以局部在线控制协调转发压力和能量支出；以序列策略利用多来源证据之间的补充关系；以接收状态估计支持发送、等待与停止的动态取舍。相应评价既观察任务或网络效果，也观察为获得这些效果付出的通信成本。

从方法上看，三个问题具有不同的可观测条件。路由能够直接利用本地队列、电量和邻居广播，但难以准确预知未来连接；内容选择可以获得候选特征，却难以用简单解析函数描述组合后的任务价值；交付阶段掌握发送历史与部分确认，真实接收集合仍然隐藏。针对这些差异分别使用在线优化、策略学习与概率状态控制，有助于把计算量放在各层最主要的不确定性上，而不要求一个模型同时学习所有通信环节。

从应用上看，无人机网络需要在有限资源下连续执行任务。关键中继提前退出、发送冗余语义或反复补发已收消息，都可能减少后续服务机会。本文通过统一观察收益与代价，为系统选择合适的路径、内容及交付动作提供依据。方法可以用于分析不同资源条件下的通信取舍，并为后续实际部署中的链路预测、消息调度和反馈机制设计提供参考。

实验采用与问题层次相对应的环境：网络仿真观察移动与负载条件下的路由表现，冻结语义轨迹用于比较共同证据基础上的内容选择，事件回放进一步引入数据失败、确认延迟和节点电量扣减。这样的组织便于解释收益来自哪个决策环节，也便于在后续系统集成时确定需要交换的状态和需要补充的实际开销。

\section{国内外研究现状}
\label{intro:related}

无人机协同通信受到拓扑、能量、消息内容和反馈状态的共同影响。现有研究分别从路由、在线资源调度、协同感知、语义表示和可靠传输等角度展开。其中，路由研究关注可用路径与传输代价，协同感知研究关注信息交换对判断结果的作用，可靠传输研究则进一步处理接收结果不确定时的后续动作。本节沿这些方向梳理已有方法，分析其解决的问题及适用条件。

\subsection{无人机自组网路由与能量控制}

飞行自组网中的节点既执行感知任务，也承担数据中继。Bekmezci等指出，三维空间中的快速运动、频繁变化的连接关系与受限机载资源，是飞行自组网区别于传统移动网络的基本因素\cite{bekmezci2013fanet}。Zeng等分析了无人机通信的体系结构和信道特征\cite{zeng2016uav}；Mozaffari等进一步从空中基站与蜂窝终端两类角色讨论部署、轨迹和资源分配\cite{mozaffari2019tutorial}。无人机的灵活部署能够改善覆盖，但同时要求协议适应运动和能量变化。

传统路由主要依赖拓扑或地理位置确定下一跳。AODV按需发现路由，利用目的序列号判断路径的新旧程度，减少无业务时的维护开销\cite{aodv}；OLSR主动维护拓扑，通过多点中继减少控制消息的重复转发\cite{olsr}。两者在路径发现时机与维护开销之间采用不同取舍，节点快速移动时，已有路径都需要随连接关系更新。GPSR根据邻居与目的节点的位置贪婪转发，在局部无法前进时切换到周边转发\cite{gpsr}。地理路由利用局部信息降低了对完整拓扑的依赖，但位置进度不能直接表示传输成功率与实际能耗。这些方法具有明确的转发规则，却可能使流量长期集中于位置有利的少数节点，路径可用性并不能保证能量均衡。

针对节点高速移动，后续方法引入链路稳定性、运动趋势和剩余能量。Lakew等归纳了拓扑变化、维护开销与网络寿命之间的制约关系\cite{intro4}。QGeo利用Q学习修正地理路由中的下一跳价值\cite{qgeo}，QTAR进一步将拓扑状态纳入学习过程\cite{qtar}，从而根据交互结果调整转发选择。与固定规则相比，学习方法可以把历史通信结果用于候选评价，但其优化方向仍取决于状态和回报中是否包含相应的资源因素。Li等提出节能快速路由方法，结合电池、信道和邻居状态联合选择下一跳与发射功率\cite{drlfr}，使单次传输的收益与实际代价共同参与决策。

路由与功率存在耦合：提高功率可能改善成功率，也会增加能耗与干扰；降低功率虽然节省单次支出，却可能产生更多重传。连续任务还要求考虑历史消耗，避免当前收益较高的路径因反复使用而提前失效。本文第一项研究由此引入长期能量反馈，在局部路由中协调转发效率与节点持续服务能力。

\subsection{基于队列稳定性的在线资源调度}

业务到达和无线信道难以事先准确预测，依赖完整未来状态的资源分配不易用于持续运行系统。Tassiulas和Ephremides建立了受约束排队系统的稳定性分析基础，利用队列积压表示尚未满足的服务需求\cite{tassiulas1992stability}。队列随业务到达增加、随实际服务减少，因而能够反映当前负载与服务能力之间的差额。Neely系统研究了李雅普诺夫漂移及漂移加惩罚方法，将长期随机约束转化为依赖当前队列状态的逐时隙决策\cite{neely2010stochastic}。对于平均能量等约束，可构造虚拟队列累计支出偏差，使历史资源压力持续影响后续动作。虚拟队列不代表实际缓存，而是把多次决策造成的资源偏差保留在状态中；当消耗持续超过目标时，后续控制会提高对能量代价的重视程度。

反压路由根据邻居之间的队列差调节流量，能够响应负载变化，但可能因绕行增加时延。Ying等结合最短路径和反压路由，以缩短转发路径\cite{related2}；Yu和Neely研究联合速率控制与路由，分析效用最优性差距与有限队列长度的关系\cite{algo1}。这些工作说明，队列稳定性需要与业务时延和路径代价共同考虑。在飞行自组网中，负载较低的邻居还可能具有较差链路或较少剩余能量，仅按队列差分流不能覆盖全部转发条件，需要与可行链路集合及通信代价结合。

在线队列方法也被用于无人机资源管理。Lin等面向太阳能无人机边缘计算，联合优化资源分配和三维轨迹，在队列与能量约束下组织感知、卸载和通信\cite{intro2}。不过，能量采集、轨迹可控性以及全局状态的可获得程度，都会改变时隙子问题。面向可充电系统的平均收支约束，与固定电池下任务结束前的可用能量约束也不能直接等同。本文主要利用虚拟队列保留资源消耗的时间关联：LEADR在连续时隙中调节路由，C-BFC在有限任务序列中调节交付代价，两者的约束性质分别在对应模型下讨论。

\subsection{多智能体选择性通信与协同感知}

多智能体通信用于弥补单个节点观察范围的不足。RIAL和DIAL将消息与任务动作共同学习，DIAL在训练时允许梯度通过通信连接传播\cite{rial-dial}；CommNet使用连续消息连接智能体隐状态\cite{commnet}。这类方法能够从任务反馈中学习通信内容，但频繁的全体交换会使通信量随节点规模增长。

后续研究因而显式选择通信对象和时机。IC3Net以门控决定是否发送\cite{ic3net}，SchedNet在共享介质和并发发送限制下学习广播调度\cite{schednet}，TarMAC则学习向哪些接收者传递消息\cite{tarmac}。三者分别对应是否通信、由谁通信和向谁通信，使资源分配成为可学习的动作。其共同特点是把通信选择嵌入任务策略，依据协作收益训练消息或门控。当消息具有不同长度、可靠性和表达粒度时，是否发送还需与发送成本一同评价，仅限制同时发送的节点数量不足以描述实际资源支出。

在协同感知中，观察差异通常具有空间含义。V2VNet交换压缩中间特征，用于联合感知与预测\cite{v2vnet}；DiscoNet借助知识蒸馏学习协作图，对不同区域的邻居特征赋予不同权重\cite{disconet}。中间特征可减少原始观测传输，但仍需处理视野重叠造成的冗余。Who2com通过学习握手匹配请求方和信息提供方\cite{who2com}，When2com共同学习协作关系与通信时机\cite{when2com}；Where2comm利用置信图筛选稀疏空间特征\cite{where2comm}，How2comm进一步结合空间与通道选择、时延补偿和时空融合\cite{how2comm}。前两者侧重选择协作节点及通信关系，后两者将选择范围细化到空间区域和特征通道，能够更有针对性地保留互补信息。它们也说明，信息价值与接收者的已有观察密切相关，单独提高发送方的局部置信度并不等于提高协作收益。

按空间位置筛选信息适合具有区域对应关系的特征融合。通信对象若转为任务标签和摘要，其作用还取决于表达中保留了哪些判断依据，同样长度的两条消息也可能覆盖不同语义。因而，空间稀疏性能够指导减少冗余，但还需要结合任务输出与消息粒度评价传输收益。

消息选择还具有顺序依赖：一次融合可能改变后续证据的作用，仅依赖静态置信度难以判断何时停止。空间上的互补关系能够帮助发现缺失区域，但不能完全说明一条新消息能否修正当前判断，尤其当候选包含不同粒度的标签和摘要时，重复语义与不同来源的一致支持还需要区别对待。图注意力网络通过关系权重表示候选之间的联系\cite{gat}，强化学习则提供了根据长期回报学习选择序列的工具\cite{sutton2018rl}。DQN估计离散动作价值，Double DQN分离选择与评估以缓解过高估计\cite{mnih2015dqn,vanhasselt2016double}；PPO以裁剪目标抑制过大的策略更新，结合广义优势估计改善训练稳定性\cite{ppo,gae}。这些工具提供表示和求解方式，仍需根据具体通信任务确定状态、动作与回报。本文第二项研究结合图结构和通信历史，在候选、中心需求与剩余预算的共同状态下学习证据选择。

\subsection{任务语义通信与视觉语言证据}

经典信息论用熵、信道容量和率失真描述有限资源下的信息恢复能力\cite{shannon1948,cover2006information}。协同判断通常不要求重建全部观测，而需要保留支持决策的信息。信息瓶颈通过约束中间表示的冗余并保留目标相关性，为任务压缩提供了分析思路\cite{tishby2000ib}。

深度联合信源信道编码把编码器、信道与解码器联合训练。Bourtsoulatze等学习图像到信道符号的映射，以改善受限带宽或信噪比下的传输\cite{bourtsoulatze2019deepjscc}；DeepSC利用Transformer学习文本语义表示与传输映射\cite{xie2021deepsc,vaswani2017attention}。Shao等面向边缘推理，以信息瓶颈保留有效特征，并使用可变长度表示适应信道\cite{shao2021task}。G\"und\"uz等强调应结合上下文和接收端功能评价通信效果\cite{semantic-communications-jsac}，使优化目标从恢复精度延伸到任务完成质量。这些方法主要研究给定内容如何形成适合任务和信道的表示；多源协同还需要决定哪些节点的内容应被选入通信过程。表示压缩与来源选择因此是相互联系的两个问题，前者节省单条消息的资源，后者减少缺少新增作用的消息。

视觉语言模型使任务信息能够以显式证据组织。CLIP通过图文对比学习建立共享表示\cite{clip}；BLIP-2以查询Transformer连接冻结视觉编码器和语言模型\cite{blip2}；LLaVA通过视觉指令微调支持围绕图像的回答与解释\cite{llava}。这些能力可用于提取对象、属性和局部判断，减少多节点系统对完整图像交换的依赖。

但同一图像面对不同问题时，所需信息并不相同。Shih等研究视觉问答中的问题相关区域选择\cite{where-to-look-vqa}；VQA v2利用同一问题对应不同答案的图像对减少语言先验偏差\cite{vqa-v2}；GQA以场景图构造组合推理问题\cite{gqa}。这些研究提示，局部观察是否有用需要相对于问题判断，能够生成完整描述并不意味着其中每部分都值得传输。COCO的图像和标注则为局部观测构造提供基础\cite{coco}。当多个节点覆盖同一对象时，还需要保留证据来源与观察范围，才能分析它们是相互补充还是重复表达。本文据此设置任务相关证据层，保留发送者、观察范围及不同消息粒度，在既定语义模型上学习调度。该分工便于分析通信策略本身的作用，也为后续区分新证据和重复发送提供统一对象。

\subsection{不完备反馈下的可靠传输与有限前瞻决策}

消息选定后仍需经过实际链路。自动重传请求根据确认决定是否再次发送，停等、回退和选择重传分别在控制复杂度、缓存与效率之间取舍\cite{lin1984arq}。传统机制以分组接收为主要依据，而语义任务还需要考虑重传能否改善判断：一份正确到达的分组可能作用有限，一份缺失的关键证据却可能决定结论。

已有工作将任务目标引入重传。Sheng等在协同感知中利用重要性图提取特征，结合语义错误检测和HARQ提高衰落信道下的可靠性\cite{sheng2024harq}。Li等研究语义HARQ，在重复语义信号、多轮接收融合和增加语义冗余之间设计不同方案\cite{li2025semharq}。这些方法将重传与感知质量联系起来，使重传内容及多轮融合不再仅围绕比特错误组织。不过，改进一次重传的有效性与判断当前是否需要重传仍有区别。在确认延迟或丢失时，发送端还需估计接收状态，并比较重传已有证据与发送其他证据的机会成本。

部分可观测马尔可夫决策过程以信念分布表示隐藏状态，使动作能够利用不完整反馈和历史记录\cite{kaelbling1998pomdp}。Huang等以信息效用损失刻画收发状态不一致的影响，在卫星融合网络中联合选择采样、编码分组与重传\cite{uoi2024satellite}。这类状态估计方法与语义重传的联系在于：发送收益既由消息内容决定，也由接收方可能已经获得哪些内容决定。对于无人机语义交付，未收到确认可能意味着数据失败，也可能只是确认未到，两种情况需要不同的动作。立即补发可能增加重复传输，持续等待则会占用有限任务窗口，因此反馈不确定性本身需要进入决策状态。

完整信念规划会随候选和反馈分支增加而迅速变大。模型预测控制采用有限时域规划并随新状态滚动更新，为约束下的在线计算提供思路\cite{mayne2000mpc}；离散动作学习也可估计后续收益，SD-SAC针对离散软演员评论家的价值低估和训练不稳定进行改进\cite{zhou2024sdsac}。本文第三项研究采用显式接收信念与有限前瞻，把新发、重传、等待和停止放在同一过程内比较。等待收益由确认能否及时返回、能否改变后续选择共同决定，因此需要随任务状态计算。

\subsection{现有工作的不足与本文研究定位}

在动态路由与资源调度方面，现有方法能够利用位置、链路或历史回报改善传输，但单次代价较低的路径可能被反复使用，关键节点仍会较快消耗电量。队列方法提供了保留长期资源压力的途径，其具体约束与可观测状态仍需结合任务确定。固定通信预算下，如何在局部信息和有限计算量下协调业务积压、链路功耗及节点能量分布，是持续执行任务时需要解决的问题。

在语义通信与协同感知方面，选择来源、空间区域和特征通道能够减少无关信息，但候选作用还会随接收方已有信息变化。任务证据具有不同来源和表达粒度，静态置信度或固定发送数量难以同时反映其可靠性、补充作用及链路成本。因此，需要在任务状态和通信历史下评价候选序列，并根据后续增益决定停止时机。

在语义交付方面，重要性筛选与多轮融合改善了重传内容的利用，但接收状态未知时仍需判断是否继续投入资源。数据失败和确认未到会产生相同的本地表象，却对应不同的重传需求。显式状态估计能够保留这种差异，有限前瞻则需要进一步权衡反馈带来的信息、等待占用的窗口以及后续任务的能量需求。

据此，本文从通信承载、内容选择和交付控制三个层次展开研究，分别处理传输机会的持续性、任务证据的有效性与接收结果的不确定性，在相应模型和实验中分析有限资源下的协同通信取舍。


\section{本文研究内容及章节安排}
\subsection{主要研究内容}

围绕上述三个问题，本文形成以下研究内容。

（1）高动态无人机网络中的能量均衡路由与功率控制。建立包含节点移动、无线链路、数据队列和能量动态的系统模型，采用李雅普诺夫漂移加惩罚构造时隙目标，以虚拟能量队列描述相对于平均额度的支出压力。在固定当前干扰的条件下求解候选链路功率，再结合地理进度与邻居剩余能量选择下一跳，利用相对运动预测调整邻居更新。由此形成仅依赖本地与邻居信息的 LEADR 分布式算法。理论分析说明累计预算关系及功率子问题的求解条件，实验从节点规模、通信范围和飞行速度考察吞吐、丢包、能量分布与网络寿命。

（2）面向协同视觉推理的链路感知选择性语义通信。利用冻结视觉语言模型生成结构化轨迹，将视觉问题映射为五级证据可见性判别，使不同通信策略共享局部证据与参考判断。EviComRL 首先依据语义需求和链路成本检索候选，再联合编码发送方、中心信念、消息粒度、可靠性与已选历史。通过代理预热及带动作掩码的近端策略优化学习消息序列和停止时机，显式处理预算、接收容量与送达门限。实验比较基础策略、协同感知通信原则的轨迹适配和网络结构对照，分析资源、覆盖、链路及关键组件的影响。

（3）不完备反馈下的语义交付与能量协同控制。在候选内容给定的基础上，建立随机数据送达、累计确认及跨任务电量更新模型。C-BFC 以轻量网络估计候选集合效用，对同一任务中输出同类别的非空接收集合进行一致性修正；根据最新累计确认与其后的发送历史递推接收信念，采用两步发送前瞻与反馈等待价值比较可行动作。虚拟能量队列和剩余额度共同调节发送价格，使任务内交付与跨任务资源支出保持联系。实验改变信道波动、反馈延迟、任务时限和能量条件，并通过消融分析效用修正、状态估计与等待机制的作用。

\subsection{整体研究方案}

\begin{figure}[htbp]
\centering
\begin{tikzpicture}[font=\small,>=Stealth,
 box/.style={draw,rounded corners=2pt,align=left,text width=37mm,minimum height=44mm,inner sep=6pt,anchor=north},
 title/.style={font=\small\bfseries,align=center,text width=39mm}]
\node[box] (a) at (0,2.5) {\textbf{LEADR}\\高动态多跳网络\\数据队列与能量压力\\下一跳和发送功率\\吞吐、丢包与能量分布};
\node[box] (b) at (4.8,2.5) {\textbf{EviComRL}\\任务相关局部证据\\语义状态与通信预算\\发送来源和消息粒度\\推理质量与通信成本};
\node[box] (c) at (9.6,2.5) {\textbf{C-BFC}\\候选证据的实际交付\\接收信念与剩余能量\\发送、等待和结束\\完成率与交付能耗};
\node[title,above=3mm of a.north] {研究一：持续传输};
\node[title,above=3mm of b.north] {研究二：内容选择};
\node[title,above=3mm of c.north] {研究三：可靠交付};
\draw[->,thick] (a.east) -- (b.west);
\draw[->,thick] (b.east) -- (c.west);
\draw[->,dashed] (a.south) -- ++(0,-.85) -| node[below,pos=.25,font=\footnotesize]{能量状态与在线队列控制思想} (c.south);
\end{tikzpicture}

\bicaption{本文三项研究的对象、方法与衔接关系}{Research objects, methods, and connections among the three studies}
\label{fig:intro:chain}
\end{figure}

图~\ref{fig:intro:chain} 展示了三项工作的关系。首先，从持续业务流出发描述节点、连接、队列和通信能量，研究如何在线分配多跳传输机会。其次，在任务语义表示基础上，把这些机会与候选证据的价值联系起来，研究通信前的内容筛选。最后，把已筛选内容放入具有失败和反馈事件的传输过程，研究实际交付以及任务之间的能量分配。

两项语义研究的衔接体现在候选与接收状态的进一步细化。内容选择需要在较多来源和粒度组合间比较，适合通过共享特征及策略网络表示候选关系。

交付控制则在较小集合上反复处理尝试与确认，可以显式维护接收组合的概率。前者决定潜在语义收益来自哪些内容，后者决定在何种接收状态下继续投入资源，两者采用不同的计算方式以适应各自决策规模。

从状态与动作看，路由以本地数据积压、邻居链路和能量状态为输入，输出下一跳与功率；内容选择以候选证据、中心信念和通信历史为输入，输出发送方、粒度或停止动作；交付控制以给定候选、发送记录和累计确认为输入，输出继续发送、补发、等待或结束任务的决定。三项研究采用相应的网络仿真、语义轨迹和交付事件回放，其共同依据是资源约束下的有效服务，具体指标则随研究对象确定。这样既能解释每项方法改变了哪个环节，也能保留后续联合实现所需的输入输出关系。

能量状态贯穿整个方案。第一项工作通过队列将历史支出反馈到当前路由，第三项在有限任务序列中延续这一思想，调节不同发送节点的资源价格。两者分别受到网络时隙与任务事件约束，共享资源调节思路，各自推导相应性质。本文按研究环节分别验证方法，完整系统还需要把路由输出的速率与排队时延、上层候选以及交付记录接入共同运行过程。

\subsection{章节安排}

全文共六章。第~\ref{ch:introduction} 章介绍研究背景，提出三个问题，梳理相关工作并说明整体方案。第~\ref{ch:foundations} 章统一网络、资源、语义证据和反馈概念，介绍后续使用的在线优化与学习控制基础。

第~\ref{rq1:chapter} 章给出 LEADR 的系统模型、漂移分析、路由功率算法及网络实验；第~\ref{rq2:chapter} 章说明 EviComRL 的轨迹构造、策略网络和训练过程，评价多种通信条件下的选择效果；第~\ref{rq3:chapter} 章建立 C-BFC 的交付模型，展开效用一致性、接收信念和反馈控制，并进行多条件及组件实验。第~\ref{ch:conclusion} 章总结主要结论，讨论联合设计、任务扩展与实际部署方向。
